\subsection{Seleção das Fontes de Dados}

A seleção das fontes de dados foi conduzida por meio de uma busca exploratória em mecanismos de pesquisa, repositórios públicos e documentações oficiais relacionadas à segurança cibernética e vulnerabilidades de software. O objetivo foi identificar bases públicas capazes de fornecer informações complementares sobre identificação de vulnerabilidades, severidade, exploração conhecida, fraquezas de software e pacotes afetados, passíveis de integração por meio de identificadores compartilhados.

As fontes candidatas foram avaliadas quanto à disponibilidade pública dos dados, documentação técnica, forma de acesso (APIs, feeds, arquivos ou repositórios Git), frequência de atualização, granularidade das informações, restrições de licenciamento e potencial de reutilização em pesquisas acadêmicas. Foram consideradas apenas bases que disponibilizassem acesso aberto aos metadados e permitissem integração com outras fontes utilizando identificadores padronizados, principalmente o CVE-ID.

Ao final do processo, foram identificadas inicialmente 30 bases públicas, das quais 10 foram analisadas tecnicamente e 7 selecionadas para análise comparativa. A versão atual do datalake implementa quatro dessas fontes: CISA Known Exploited Vulnerabilities (KEV), CVE List V5, National Vulnerability Database (NVD) e GitHub Advisory Database. Além dessas, o modelo de dados prevê a futura incorporação de bases complementares, como CWE, OSV, EPSS, Exploit-DB e outras fontes de threat intelligence.

As principais fontes utilizadas na implementação são apresentadas na Tabela~\ref{tab:data_sources}.

\begin{table}[htbp]
    \centering
    \caption{Fontes implementadas na versão atual do dataset}
    \label{tab:data_sources}
    \small
    \renewcommand{\arraystretch}{1.15}
    \begin{tabular}{|>{\raggedright\arraybackslash}p{3.2cm}|>{\raggedright\arraybackslash}p{2.8cm}|>{\raggedright\arraybackslash}p{6.5cm}|}
        \hline
        \textbf{Fonte}           &
        \textbf{Forma de acesso} &
        \textbf{Principais informações}                                                \\
        \hline

        CISA KEV                 &
        JSON Feed                &
        Exploração conhecida, data de inclusão, ação recomendada e prazo de mitigação. \\

        \hline

        CVE List V5              &
        Repositório Git          &
        Registro canônico de CVEs, descrições, produtos afetados e referências.        \\

        \hline

        NVD                      &
        REST API                 &
        CVSS, CWE, CPE, métricas de severidade e metadados técnicos.                   \\

        \hline

        GitHub Advisory Database &
        Repositório Git          &
        Pacotes afetados, ecossistemas, versões vulneráveis/corrigidas e advisories.   \\

        \hline
    \end{tabular}
\end{table}
